\usepackage{langsci-optional}
\usepackage{langsci-gb4e}
\usepackage{langsci-lgr}

% Load xeCJK for Japanese text (needed for マクドナルド etc)
\usepackage{xeCJK}
% Try common CJK serif fonts to match book's serif style
\IfFontExistsTF{Noto Serif CJK JP}{
  \setCJKmainfont{Noto Serif CJK JP}
  \setCJKsansfont{Noto Sans CJK JP}
  \setCJKmonofont{Noto Sans Mono CJK JP}
}{
  \IfFontExistsTF{Hiragino Mincho ProN}{
    \setCJKmainfont{Hiragino Mincho ProN}
    \setCJKsansfont{Hiragino Sans GB}
    \setCJKmonofont{Hiragino Sans GB}
  }{
    % Fallback - still prefer serif-style for main
    \setCJKmainfont{Arial Unicode MS}
    \setCJKsansfont{Arial Unicode MS}
    \setCJKmonofont{Arial Unicode MS}
  }
}
\usepackage{fontspec}
% Setting up Hindi font
\newfontfamily\hindifont{Noto Serif Devanagari}  % Replace [HindiFont] with your Hindi font

\usepackage{mdframed}
\usepackage{tcolorbox}
\tcbuselibrary{breakable}
\usepackage{tipa}
% contour and soul removed (no longer needed after underline cleanup)
\usepackage{ulem}  % retained for \xout{} in Ch 15 (ellipsis strikethrough)
\usepackage{langsci-tbls}
% Override tblsfilled: upstream uses frame-based fill that doesn't render
% reliably; use colback instead, and add breakable for page-spanning boxes
\RenewTColorBox{tblsfilled}{m O{black!12}}
  {
    enhanced,
    breakable,
    title = {#1},
    colback = #2,
    colframe = #2,
    boxrule = 0pt,
    boxsep = 0pt,
    toptitle = 5mm,
    top = 5mm,
    bottom = 5mm,
    left = 5mm,
    right = 5mm,
    sharp corners = all,
    coltitle = black,
    fonttitle = \normalfont\sffamily\bfseries\large,
  }
\usepackage{enumitem}
\usepackage{multicol}
\usepackage{forest}
\usepackage{graphicx}
\usepackage{wrapfig}
\usetikzlibrary{bending}
\usepackage{subcaption}
\usepackage{changepage}

% Fix header height warnings
\usepackage{geometry}
\geometry{head=27.2pt}

%\usepackage{fontspec}

%\setmainfont{Times New Roman} % Your main document font
%\newfontfamily\japanesefont{IPAMincho} % Japanese font
%\newfontfamily\hindifont[Script=Devanagari]{Lohit Devanagari} % Hindi font


\usepackage{listings}
\lstset{basicstyle=\ttfamily,tabsize=2,breaklines=true}

\usepackage{tikz}
\usetikzlibrary{tikzmark, positioning}

\usepackage{longtable}

\usepackage{dialogue}

\usepackage{paracol}

% mdframed already loaded above (line 28)

\usepackage{pifont}

\newfontfamily\ethiopic[Script=Ethiopic]{Abyssinica SIL}
\newfontfamily\persian[Script=Arabic]{Amiri}

% ==========================================================
%                  INDEX SETUP (REVISED)
% ==========================================================
% This setup uses a simpler definition to avoid conflicts.
% PLEASE REPLACE ALL OTHER INDEXING CODE WITH THIS BLOCK.

% 1. Load the package for indexing
\usepackage{imakeidx}

% 2. Create the three separate indexes + dummy output index
\makeindex[name=output, title=Output Index, columns=2]
\makeindex[name=sbj, title=Subject Index, columns=2, intoc]
\makeindex[name=lan, title=Language Index, columns=2, intoc]
\makeindex[name=aut, title=Author Index, columns=2, intoc]

% 3. Define simple, robust custom commands
% Using \renewcommand avoids conflicts with the langscibook class.
\renewcommand{\is}[1]{\index[sbj]{#1}}
\renewcommand{\il}[1]{\index[lan]{#1}}
\renewcommand{\ia}[1]{\index[aut]{#1}}
% ==========================================================

% csquotes for \enquote{} (smart quotation marks)
\usepackage{csquotes}

% Semantic typography macros
\newcommand{\mention}[1]{\textit{#1}}        % metalinguistic mention (word as object)
\newcommand{\mentionhead}[1]{⟨#1⟩}           % mention in headings (angle brackets)
\providecommand{\term}[1]{\textsc{#1}}       % technical term (small caps)